% fig_posterior_zoom.tex
%
% From the deterministic map to the posterior of one plant and one link.
%   a  deterministic map (Box 2 key)
%   b  average posterior per plant and category (heatmap)
%   c  the chosen plant: expected links per category
%   d  the chosen link's category map, camera records as asterisks
%
% The panels and anchors.tex are written by bayesian_empirical.R, Section 13,
% into posterior_zoom/. anchors.tex says where the chosen plant row and link
% cell sit inside each panel, so the boxes and connectors below follow the
% panels whenever the script is rerun (another plant, link or run).
%
% Compile from this folder:        pdflatex fig_posterior_zoom.tex
% Place in the manuscript:         \includegraphics[width=\linewidth]{fig_posterior_zoom.pdf}
% The panels are drawn at their saved size, and the figure is scaled as a whole.

\documentclass[tikz, border = 6pt]{standalone}
\usepackage{graphicx}
\usepackage[scaled = 0.92]{helvet}
\renewcommand{\familydefault}{\sfdefault}
\usetikzlibrary{calc, arrows.meta}

\newcommand{\fzdir}{posterior_zoom}
% written by bayesian_empirical.R, Section 13: Helichrysum stoechas x Meliscaeva auricollis, B-uniform
\newcommand{\fzaL}{0.0654}
\newcommand{\fzaB}{0.4428}
\newcommand{\fzaR}{0.9854}
\newcommand{\fzaT}{0.9852}
\newcommand{\fzbL}{0.3151}
\newcommand{\fzbB}{0.0231}
\newcommand{\fzbR}{0.8047}
\newcommand{\fzbT}{0.7537}
\newcommand{\fzcL}{0.3288}
\newcommand{\fzcB}{0.1452}
\newcommand{\fzcR}{0.9753}
\newcommand{\fzcT}{0.8890}
\newcommand{\fzdL}{0.0187}
\newcommand{\fzdB}{0.7453}
\newcommand{\fzdR}{0.7292}
\newcommand{\fzdT}{0.9447}
\newcommand{\fzaRowB}{0.5445}
\newcommand{\fzaRowT}{0.5929}
\newcommand{\fzaColL}{0.3860}
\newcommand{\fzaColR}{0.4060}
\newcommand{\fzdRowB}{0.7619}
\newcommand{\fzdRowT}{0.9280}
\newcommand{\fzdColL}{0.2663}
\newcommand{\fzdColR}{0.2817}
\newcommand{\fzbRowB}{0.1601}
\newcommand{\fzbRowT}{0.2253}


\definecolor{fzink}{HTML}{2B2B3A}
\definecolor{fzband}{HTML}{7B7BB8}

% \fzat{name}{panel}{fx}{fy}
% A named point inside a panel node, fx and fy as fractions of the panel
% file's width and height, measured from its bottom left corner.
\newcommand{\fzat}[4]{%
  \path let \p1 = (#2.south west), \p2 = (#2.north east) in
    coordinate (#1) at ({\x1 + (#3)*(\x2 - \x1)}, {\y1 + (#4)*(\y2 - \y1)});}

\tikzset{
  panel/.style  = {inner sep = 0pt, outer sep = 0pt},
  letter/.style = {font = \bfseries\Large, anchor = south east, inner sep = 1pt},
  focus/.style  = {draw = fzink, line width = 0.9pt},
  cell/.style   = {draw = fzink, line width = 1.3pt},
  band/.style   = {fill = fzband, fill opacity = 0.16},
  edge/.style   = {draw = fzband, line width = 0.6pt},
  halo/.style   = {draw = white, line width = 2.6pt},
  link/.style   = {draw = fzink, line width = 0.9pt, -{Latex[length = 2.4mm]}},
}

\begin{document}
\begin{tikzpicture}

  % ---- panels --------------------------------------------------------------
  % b and c along the top, d under b, a to the left. Every panel's top edge is
  % flush with b's, and a is as tall as b, the gap and d together.
  \node[panel, anchor = north west] (b) at (0, 0)
    {\includegraphics{\fzdir/panel_b.pdf}};
  \node[panel, anchor = north west] (c) at ($(b.north east) + (0.9, 0)$)
    {\includegraphics{\fzdir/panel_c.pdf}};
  \node[panel, anchor = north west] (d) at ($(b.south west) + (0, -0.9)$)
    {\includegraphics{\fzdir/panel_d.pdf}};
  \node[panel, anchor = north east] (a) at ($(b.north west) + (-1.5, 0)$)
    {\includegraphics{\fzdir/panel_a.pdf}};

  % (\pn, not \p, which the let syntax in \fzat reserves)
  \foreach \pn in {a, b, c, d}
    \node[letter] at (\pn.north west) {\pn};

  % ---- anchor points -------------------------------------------------------
  % a: the plant row, the link cell, and the row's right end
  \fzat{aRowSW}{a}{\fzaL}{\fzaRowB}
  \fzat{aRowNE}{a}{\fzaR}{\fzaRowT}
  \fzat{aRowSE}{a}{\fzaR}{\fzaRowB}
  \fzat{aRowE} {a}{\fzaR}{(\fzaRowB + \fzaRowT) / 2}
  \fzat{aCellSW}{a}{\fzaColL}{\fzaRowB}
  \fzat{aCellNE}{a}{\fzaColR}{\fzaRowT}
  \fzat{aCellE} {a}{\fzaColR}{(\fzaRowB + \fzaRowT) / 2}

  % b: the plant row, from the plant name to the last column
  \fzat{bRowSW}{b}{0.01}{\fzbRowB}
  \fzat{bRowNW}{b}{0.01}{\fzbRowT}
  \fzat{bRowNE}{b}{\fzbR}{\fzbRowT}
  \fzat{bRowSE}{b}{\fzbR}{\fzbRowB}

  % c: the panel's outer left edge, level with the plotting area, so the band
  % stops before the category labels
  \fzat{cNW}{c}{0}{\fzcT}
  \fzat{cSW}{c}{0}{\fzcB}

  % d: the plant row and the link cell
  \fzat{dRowSW} {d}{\fzdL}{\fzdRowB}
  \fzat{dRowNE} {d}{\fzdR}{\fzdRowT}
  \fzat{dRowW}  {d}{\fzdL}{(\fzdRowB + \fzdRowT) / 2}
  \fzat{dCellSW}{d}{\fzdColL}{\fzdRowB}
  \fzat{dCellNE}{d}{\fzdColR}{\fzdRowT}
  \fzat{dCellW} {d}{\fzdColL}{(\fzdRowB + \fzdRowT) / 2}

  % ---- one plant: map row -> heatmap row -> its distribution ---------------
  \fill[band] (aRowNE) to[out = 0, in = 180] (bRowNW) --
              (bRowSW) to[out = 180, in = 0] (aRowSE) -- cycle;
  \draw[edge] (aRowNE) to[out = 0, in = 180] (bRowNW);
  \draw[edge] (aRowSE) to[out = 0, in = 180] (bRowSW);

  \fill[band] (bRowNE) to[out = 0, in = 180] (cNW) --
              (cSW) to[out = 180, in = 0] (bRowSE) -- cycle;
  \draw[edge] (bRowNE) to[out = 0, in = 180] (cNW);
  \draw[edge] (bRowSE) to[out = 0, in = 180] (cSW);

  \draw[focus] (aRowSW) rectangle (aRowNE);
  \draw[focus] (bRowSW) rectangle (bRowNE);

  % ---- one link: map cell -> the plant's row of the category map -----------
  % along the plant row to the map's edge, then down into the strip, which is
  % that same row lifted out of the map
  \draw[focus] (dRowSW) rectangle (dRowNE);
  \draw[halo] (aCellE) -- (aRowE) to[out = 0, in = 180] (dRowW) -- (dCellW);
  \draw[link] (aCellE) -- (aRowE) to[out = 0, in = 180] (dRowW) -- (dCellW);

  \draw[cell] (aCellSW) rectangle (aCellNE);
  \draw[cell] (dCellSW) rectangle (dCellNE);

\end{tikzpicture}
\end{document}
